\chapter{Preface}

\section{Project Description}

Assembly Copilot is a computer-vision system that observes an operator's
egocentric video while an assembly procedure is being performed. It estimates
which procedure step is currently in progress, announces completed steps, and
surfaces out-of-order activity for review. The project connects procedure-step
recognition research with a working demonstration that can be replayed in a
browser, served through a live video path, and displayed on a HoloLens overlay.

The demonstrator uses a ten-part turbofan engine model as the assembly object.
The model is deliberately small enough to support rapid experimentation while
retaining the practical characteristics of a production task: hand-object
occlusion, similar-looking parts, variable step duration, incomplete views, and
operators who do not follow exactly the same order.

\section{Project Value}

The intended value is an additional source of situational awareness for an
operator or supervisor. A live system can reduce the need to manually track
procedure progress, make deviations visible earlier, and provide a consistent
record for post-task review. The system is an assistive demonstrator; it is not
a replacement for approved work instructions, inspection, or human sign-off.

\section{Documentation Scope}

This document records the development path from preliminary procedure-step
recognition experiments to the current live inference demonstration. It covers
the project objective, data collection and labeling tools, model development,
streaming inference, application presentation, HoloLens integration, measured
results, limitations, and recommended next steps. Proprietary internal model
implementation details are described at the level needed to understand the
system behavior and engineering decisions.

\section{Key Outcomes}

\begin{table}[H]
\centering
\small
\begin{tabularx}{\textwidth}{@{}p{0.22\textwidth}X@{}}
\toprule
\textbf{Area} & \textbf{Documented outcome} \\
\midrule
Research & Controlled progression from a VideoMAEv2 baseline to a fused
motion-and-appearance model with the best IndustReal result. \\
Data & Own egocentric engine-model recordings, dense labels, and stress
variants for normal and non-nominal procedure behavior. \\
Deployment & Causal temporal inference with fixed-lag commitment, replay and
live services, and a HoloLens presentation path. \\
Evidence & Separate offline, online, robustness, latency, and interface
evidence so the demonstrator is not represented by one number alone. \\
\bottomrule
\end{tabularx}
\caption{Summary of the project outcomes covered in this document.}
\label{tab:key-outcomes}
\end{table}
